\subsection{Sofwarearchitektur}
Für die Umsetzung des Bewerbermanagementsystems wird eine klassische 3-Schichten-Architektur gewählt. 
Diese unterteilt die Anwendung in Präsentationsschicht, Anwendungslogik und Datenhaltung. 
Die Präsentationsschicht bildet die Benutzeroberfläche für Bewerber:innen und Recruiter:innen ab, während die Geschäftslogik in PHP umgesetzt wird und die Verarbeitung von Bewerbungen, Statuswechseln und Verwaltungsfunktionen steuert. 
Die Datenhaltung erfolgt über eine SQLite-Datenbank, in der Stellenanzeigen, Bewerbungen, Nutzer und Dokumente strukturiert gespeichert werden.
Die gewählte Architektur eignet sich besonders für den prototypischen Charakter des Systems, da sie eine klare Trennung der Verantwortlichkeiten ermöglicht und dadurch Wartbarkeit, Übersichtlichkeit und Erweiterbarkeit verbessert. Gleichzeitig bleibt die technische Umsetzung schlank und angemessen für den Umfang eines Minimum Viable Product.
